\documentclass[11pt]{article}

\usepackage[margin=1in]{geometry}
\usepackage{amsmath,amssymb,amsthm}
\usepackage{graphicx}
\usepackage[hidelinks]{hyperref}

\newtheorem{theorem}{Theorem}
\newtheorem{lemma}{Lemma}
\newtheorem{remark}{Remark}

\newcommand{\N}{\mathbb{N}}
\newcommand{\Ttree}{T}
\newcommand{\Ball}{B}
\newcommand{\conv}{\operatorname{conv}}
\newcommand{\wclosedconv}{\overline{\operatorname{conv}}}
\newcommand{\norm}[1]{\left\lVert #1 \right\rVert}
\newcommand{\vertiii}[1]{\left|\!\left|\!\left| #1 \right|\!\right|\!\right|}
\newcommand{\abs}[1]{\left\lvert #1 \right\rvert}
\newcommand{\SCD}{\operatorname{SCD}}

\title{Improving the Unconditional Constant in Loo--Perreau's Non-SCD Renorming}
\author{Automated Math Discovery Smoke Test}
\date{June 8, 2026}

\begin{document}
\maketitle

\section*{Source and Open Question}

This packet concerns Question 5.4 from:

\begin{quote}
Marcus Loo and Yoel Perreau,
\emph{A Banach space with an unconditional basis which is not slicely
countably determined}, arXiv:2603.12947.
\end{quote}

The question appears in Section 5, page 18 of the source paper:

\begin{center}
\includegraphics[width=\textwidth]{figures/open_problem_crop.png}
\end{center}

The source paper constructs, from the binary tree space $X=X_T$, an
equivalent norm with a $2$-unconditional basis whose unit ball is not SCD.
Question 5.4 asks whether the constant can be improved to any
$k\in(1,2)$.

\section*{Imported Notation and Facts}

We use the notation of the source paper. Thus $X=X_T$ is the binary tree
space, $(e_t)_{t\in T}$ is its canonical $1$-unconditional basis, $\Ball_X$
is its original unit ball, and
\[
  \Ball_X^+=\{x\in \Ball_X:x(t)\ge 0\ \text{for all }t\in T\},\qquad
  \Ball_X^-=-\Ball_X^+.
\]
The paper defines $\Sigma$ and $\Pi$ inside $\Ball_X$ and writes
$\Omega=\Sigma\setminus\Pi$. We need two facts proved there:

\begin{enumerate}
\item $\Omega^+$ is non-empty and contained in $S_X^+$.
\item No point $x\in\Omega^+$ is an SCD point of $\Ball_X^+$.
This is Proposition 3.9 of the source paper.
\end{enumerate}

We also use the standard localized form of the SCD-point criterion: for a
bounded convex set, countable determining families of non-empty relatively
weakly open subsets are equivalent to countable determining families of
slices. This is the pointwise version of the Bourgain-lemma reduction used
in the source paper for SCD sets.

\section*{Main Claim}

\begin{theorem}
For every $k\in(1,2)$ there is a Banach space with a $k$-unconditional
basis whose unit ball is not slicely countably determined.
\end{theorem}

\begin{proof}
Fix $t\in(1/2,1)$ and set
\[
  C_t=\wclosedconv\bigl(\Ball_X^+\cup \Ball_X^- \cup t\Ball_X\bigr).
\]
The set $C_t$ is closed, convex, symmetric, and absorbing; hence it is the
unit ball of an equivalent norm, which we denote by $\vertiii{\cdot}_t$.
Moreover
\[
  t\Ball_X\subseteq C_t\subseteq \Ball_X.
\]
Consequently
\[
  \norm{x}\le \vertiii{x}_t\le t^{-1}\norm{x}\qquad (x\in X).
\]

We first estimate the unconditional constant of the canonical basis in
$(X,\vertiii{\cdot}_t)$. Let $D_\varepsilon$ be any sign-change operator,
$D_\varepsilon(\sum a_t e_t)=\sum \varepsilon_t a_t e_t$, where
$\varepsilon_t\in\{-1,1\}$. Since the original basis of $X$ is
$1$-unconditional and $C_t\subseteq\Ball_X$, we have
\[
  D_\varepsilon(C_t)\subseteq D_\varepsilon(\Ball_X)=\Ball_X.
\]
But $t\Ball_X\subseteq C_t$, so $\Ball_X\subseteq t^{-1}C_t$. Hence
\[
  D_\varepsilon(C_t)\subseteq t^{-1}C_t.
\]
Thus $\|D_\varepsilon\|_{(X,\vertiii{\cdot}_t)\to (X,\vertiii{\cdot}_t)}
\le t^{-1}$ for every choice of signs, so the basis is
$t^{-1}$-unconditional.

It remains to show that $C_t$ is not SCD. Choose $x\in\Omega^+$. We will
show that $x$ is not an SCD point of $C_t$.

Assume toward a contradiction that $(S_n)$ is a countable sequence of
slices of $C_t$ determining $x$. Let
\[
  M=\{n\in\N:S_n\cap \Ball_X^+\neq\varnothing\}.
\]
For $n\in M$, put $U_n=S_n\cap\Ball_X^+$. We claim that the family
$(U_n)_{n\in M}$ determines $x$ inside $\Ball_X^+$.

If $M=\varnothing$, then, arguing as below, choose
$z_n\in S_n\cap(\Ball_X^-\cup t\Ball_X)$ for every $n$. The determining
property would give
$x\in\wclosedconv\{z_n:n\in\N\}$. Applying a branch functional $\beta$
with $\beta(x)=1$ is impossible, since $\beta(z_n)\le t<1$ for all $n$.
Thus $M$ is non-empty. If $\N\setminus M=\varnothing$, the claim is
immediate from the fact that the original slices determine $x$ in $C_t$.
So assume from now on that both $M$ and $\N\setminus M$ are non-empty.

Take an arbitrary selection $y_n\in U_n$ for $n\in M$. If $n\notin M$,
then $S_n$ does not meet $\Ball_X^+$. Since
$C_t=\wclosedconv(\Ball_X^+\cup\Ball_X^-\cup t\Ball_X)$, each such slice
meets $\Ball_X^-\cup t\Ball_X$; choose $z_n$ there. Because the slices
$(S_n)$ determine $x$ in $C_t$, the point $x$ belongs to the closed convex
hull of the selected points $y_n$ and $z_n$.

Let $\varepsilon>0$. We may therefore find
\[
  y\in\conv\{y_n:n\in M\},\qquad
  z\in\conv\{z_n:n\notin M\},\qquad
  \lambda\in[0,1],
\]
such that, in the original norm of $X$,
\[
  \norm{x-\bigl(\lambda y+(1-\lambda)z\bigr)}<\varepsilon.
\]
Pick a branch $\beta$ of the binary tree for which
\[
  \sum_{s\in\beta} x(s)=1.
\]
The corresponding branch functional has norm one on $X$. Since
$y\in\Ball_X^+$, we have $\sum_{s\in\beta}y(s)\le 1$. Since
$z\in\conv(\Ball_X^-\cup t\Ball_X)$, we have
$\sum_{s\in\beta}z(s)\le t$. Applying the branch functional gives
\[
  (1-\lambda)(1-t)
  \le
  1-\lambda\sum_{s\in\beta}y(s)
    -(1-\lambda)\sum_{s\in\beta}z(s)
  \le
  \norm{x-\bigl(\lambda y+(1-\lambda)z\bigr)}
  < \varepsilon.
\]
Therefore $1-\lambda<\varepsilon/(1-t)$. Since $y,z\in\Ball_X$,
\[
  \norm{x-y}
  \le
  \norm{x-\bigl(\lambda y+(1-\lambda)z\bigr)}
  +(1-\lambda)\norm{y-z}
  <
  \varepsilon+\frac{2\varepsilon}{1-t}.
\]
As $\varepsilon$ was arbitrary,
$x\in\wclosedconv\{y_n:n\in M\}$. Thus the relatively weakly open family
$(U_n)_{n\in M}$ determines $x$ in $\Ball_X^+$.

By the localized SCD-point criterion, $x$ would be an SCD point of
$\Ball_X^+$, contradicting Proposition 3.9 of Loo--Perreau. Hence $C_t$ is
not SCD.

Finally, given $k\in(1,2)$, take $t=1/k$. Then $t\in(1/2,1)$, the canonical
basis of $(X,\vertiii{\cdot}_t)$ is $k$-unconditional, and its unit ball
$C_t$ is not SCD.
\end{proof}

\section*{Verification Notes}

The only numerical parameter check is
\[
  k\in(1,2)\quad\Longleftrightarrow\quad t=1/k\in(1/2,1),
\]
and the branch-functional estimate loses the factor $(1-t)^{-1}$, which is
finite precisely because $t<1$. The short script
\texttt{code/check\_interpolation\_constants.py} verifies these inequalities
for sample values and records the constant calculation used above.

\begin{remark}
The result does not address Question 5.5 of the source paper. The
construction still uses the same family $\Omega^+$ of non-SCD points, but
it does not show that the whole unit ball has no SCD points.
\end{remark}

\end{document}
